\documentclass[12pt,a4paper]{article}
\usepackage[a4paper,margin=1in]{geometry}
\usepackage{amsmath,amssymb,mathtools}
\usepackage{enumitem}
\usepackage{hyperref}
\usepackage{titlesec}
\usepackage{parskip}
\usepackage{booktabs}
\usepackage{tabularx}
\usepackage{array}
\usepackage{microtype}
\usepackage{xcolor}
\usepackage{fancyhdr}
\usepackage{lastpage}
\setlength{\parindent}{0pt}
\setlength{\parskip}{6pt}
\setlist[itemize]{leftmargin=1.5em,itemsep=3pt,topsep=3pt}
\setlist[enumerate]{leftmargin=1.7em,itemsep=4pt,topsep=3pt}

\definecolor{headingblue}{RGB}{31,55,75}

\titleformat{\section}
  {\normalfont\Large\bfseries\color{headingblue}}
  {\thesection}{0.8em}{}

\titleformat{\subsection}
  {\normalfont\large\bfseries\color{headingblue}}
  {\thesubsection}{0.8em}{}

\titleformat{\subsubsection}
  {\normalfont\normalsize\bfseries}
  {\thesubsubsection}{0.7em}{}

\hypersetup{
    colorlinks=true,
    linkcolor=headingblue,
    urlcolor=headingblue,
    pdftitle={AetherGrid AI: An Autonomous Logistics and Emergency-Dispatch Simulation},
    pdfauthor={Vishal Singla and Sparsh Verma}
}

\pagestyle{fancy}
\fancyhf{}
\lhead{\small AetherGrid AI}
\rhead{\small Project Proposal}
\cfoot{\small Page \thepage\ of \pageref{LastPage}}
\setlength{\headheight}{14pt}





\begin{document}





\begin{titlepage}
\centering

\vspace*{1.5cm}

{\Large\bfseries Thapar Institute of Engineering and Technology}

\vspace{0.8cm}

{\large Project Proposal for UCS503P}

\vspace{1.5cm}

{\Huge\bfseries AetherGrid AI}

\vspace{0.4cm}

{\Large\bfseries An Autonomous Logistics and\\
Emergency-Dispatch Simulation}

\vspace{1.5cm}

{\large\itshape A simulation-driven intelligent dispatch system for\\
dynamic routing, optimal vehicle assignment, and predictive\\
idle-vehicle repositioning}

\vfill

\begin{tabular}{rl}
\textbf{Presented by:} & Vishal Singla \\
\textbf{Roll No.:} & 1024240009 \\
\textbf{Presented by:} & Sparsh Verma \\
\textbf{Roll No.:} & 1024240011 \\
\textbf{Course:} & Software Engineering(UCS503P)\\
\textbf{Presented to:} & Dr. Zeelani\\

\end{tabular}

\vspace{1.5cm}

{\large August 2026}

\end{titlepage}





\section{Abstract}

AetherGrid AI is a simulation-based intelligent dispatch system
designed to study how a mixed fleet of emergency and logistics
vehicles can be managed in a changing urban environment. The system
models a city as a dynamic road network in which congestion, weather,
hazards, and road closures can continuously change the cost of
travel. It then combines pathfinding, vehicle-request assignment,
multi-stop route optimization, and reinforcement learning to make
dispatch decisions.

The project is designed as a seeded and reproducible discrete-event
simulation. Emergency and logistics requests arrive continuously,
are collected in short dispatch batches, and are assigned to suitable
vehicles. Dynamic A* is used for route execution, while Dijkstra is
used to calculate the travel-time matrix required by the assignment
stage. The Hungarian algorithm is used for one-to-one batch
assignment, while a Genetic Algorithm is used for multi-drop
logistics routes with time windows. A Q-Learning or Deep Q-Network
(DQN) based policy is used only for repositioning idle vehicles
between dispatch batches. This separation ensures that the learning
component cannot interrupt or invalidate an active emergency
dispatch.

A real-time Mission Control dashboard will display vehicles,
requests, routes, road conditions, telemetry, and simulation
statistics. The project will be evaluated against simple baseline
strategies using response time, service-level compliance, route
cost, energy consumption, runtime, and reproducibility. The primary
objective is not to claim a new algorithm, but to demonstrate the
correct and measurable integration of established algorithms into a
complete real-time simulation system.





\section{Introduction}

Urban emergency response and logistics operations require vehicles to
be assigned to requests quickly and intelligently. The difficulty is
not only finding a short route. The system must also decide which
vehicle should serve which request, account for changing road
conditions, handle several requests at the same time, and maintain
good vehicle positions for future demand.

A conventional nearest-vehicle strategy can work well in simple
situations, but it becomes less effective when demand is uneven,
roads become unavailable, or several vehicles compete for the same
requests. AetherGrid addresses this problem through a simulation
environment in which the complete dispatch process can be observed,
tested, and compared under controlled conditions.

The immediate engineering objective is to build a reproducible
simulation engine that combines a live cost-priced street graph with
three cooperating decision layers:

\begin{enumerate}
    \item dynamic pathfinding,
    \item batch-optimal vehicle assignment, and
    \item learned repositioning of idle vehicles.
\end{enumerate}

These layers are exposed through a real-time Mission Control
dashboard so that the behavior of the system can be observed during
live scenarios.





\section{Problem Statement}

Dispatching a heterogeneous fleet such as ambulances, fire engines,
drones, and delivery vans against a continuous stream of requests is
a difficult optimization problem. Streets, demand, vehicle
availability, and operating conditions can all change while the
system is running.

The major problems considered in AetherGrid are described below.

\subsection{Combinatorial Assignment Problem}

If several vehicles and several requests are available at the same
time, there may be many possible vehicle-request combinations.
Trying every possible combination is inefficient for real-time
dispatch.

A system therefore needs an assignment strategy that can find a
high-quality solution without testing every possible arrangement.

\subsection{Changing Road Conditions}

The shortest route under normal conditions may not remain the best
route when congestion, weather, hazards, or road closures occur.

Therefore, the road network must be treated as dynamic rather than
static.

\subsection{Myopic Dispatch}

A nearest-vehicle strategy considers the current request but may
ignore the effect of its decision on future requests. Assigning the
only nearby vehicle to a current request may leave an important area
without a suitable vehicle.

AetherGrid therefore separates immediate dispatch from future
positioning.

\subsection{Disruption Handling}

The system must continue to operate when roads are closed or hazards
are introduced. A closed road must never be selected by a route.

This requirement is important because the project is intended to
demonstrate not only normal-case routing but also graceful behavior
under disruption.





\section{Aim and Objectives}

\subsection{Aim}

The main aim of AetherGrid AI is to develop a reproducible urban
dispatch simulation that combines dynamic routing, optimal
vehicle-request assignment, multi-stop logistics optimization, and
predictive repositioning of idle vehicles.

\subsection{Objectives}

The project has the following objectives:

\begin{enumerate}
    \item Build a dynamic graph-based representation of an urban road
    network.

    \item Model congestion, weather, hazards, and road closures as
    changing road conditions.

    \item Generate emergency and logistics requests continuously
    using a controlled stochastic process.

    \item Calculate efficient routes using Dynamic A*, Dijkstra,
    and HPA* where appropriate.

    \item Build an ETA matrix between available vehicles and pending
    requests.

    \item Use the Hungarian algorithm to solve the batch vehicle-
    request assignment problem.

    \item Include deadline-aware penalties so that late emergency
    responses receive a higher assignment cost.

    \item Use a Genetic Algorithm to solve multi-drop logistics
    problems with time windows.

    \item Use Q-Learning or DQN to reposition only idle vehicles
    between dispatch batches.

    \item Develop a real-time Mission Control dashboard for
    visualization and scenario control.

    \item Maintain deterministic, seed-based simulation runs for
    testing and reproducibility.

    \item Compare intelligent strategies against simple greedy
    baselines using measurable performance metrics.
\end{enumerate}





\section{Proposed Solution}

AetherGrid is designed as a discrete-event simulation engine with a
web-based visualization layer.

The proposed system contains the following major components.

\subsection{Dynamic Road Network}

The city is represented as a graph containing locations and roads.
Every road has a current travel cost. The cost can change during the
simulation as traffic, weather, or hazards change.

Road closures are handled differently from ordinary hazards. A
hazard increases the cost of a road, whereas a closed road is removed
from the graph entirely.

\subsection{Pathfinding Suite}

The project uses multiple pathfinding methods for different purposes:

\begin{itemize}
    \item \textbf{Dynamic A*:} route execution using the current road
    costs.
    \item \textbf{Dijkstra:} calculation of shortest travel costs
    used for the vehicle-request ETA matrix.
    \item \textbf{HPA*:} hierarchical pathfinding for larger or
    long-haul routes where a complete fine-grained search may be
    unnecessarily expensive.
\end{itemize}

\subsection{Batch Dispatch}

Requests are collected for a short dispatch interval. The proposed
batch interval is:

\begin{equation}
    \Delta = 5\,\mathrm{s}
    \label{eq:batch}
\end{equation}

At each dispatch interval, the system calculates current travel
costs, generates the ETA matrix, and solves the vehicle-request
assignment problem.

\subsection{Logistics Route Optimization}

Emergency dispatch generally involves one immediate destination,
while logistics operations may involve several destinations.

For multi-drop logistics missions, AetherGrid uses a Genetic
Algorithm to search for routes that minimize travel cost while
respecting delivery time windows.

\subsection{Idle Vehicle Repositioning}

The reinforcement-learning component is intentionally restricted to
idle vehicles. It does not override a live dispatch.

Between batches, the RL policy may move an idle vehicle toward an
area where future demand is expected to be higher. This allows the
system to introduce a limited form of predictive behavior without
making the safety-critical dispatch layer dependent on the
convergence of a learning model.





\section{System Workflow}

The complete AetherGrid workflow is as follows:

\begin{enumerate}
    \item The simulation initializes the road network, vehicles,
    requests, environmental conditions, and random seed.

    \item Requests arrive according to the selected request
    generation model.

    \item Requests are temporarily stored until the next dispatch
    batch.

    \item The current road conditions are used to calculate the cost
    of each available road.

    \item Roads affected by closures are removed from the graph.

    \item Dijkstra or A* is used to calculate travel costs between
    vehicles and requests.

    \item The resulting values form the vehicle-request cost matrix.

    \item The Hungarian algorithm assigns vehicles to requests.

    \item If a request belongs to a multi-drop logistics mission,
    the Genetic Algorithm based VRPTW solver is used instead.

    \item Assigned vehicles execute their routes using Dynamic A*.

    \item Idle vehicles that remain unassigned may be repositioned
    by the reinforcement-learning policy.

    \item Vehicle states, routes, requests, and performance metrics
    are sent to the dashboard.

    \item After a mission is completed, the vehicle returns to the
    idle pool and becomes available for the next batch.
\end{enumerate}

This creates a continuous control loop in which the result of one
dispatch cycle affects the state of the next cycle.





\section{Mathematical Formulation}

This section defines the mathematical models used by AetherGrid.
Each equation is followed by a direct explanation of its variables
and purpose.





\subsection{Dynamic Road Network}

The urban environment is represented as a weighted graph:

\begin{equation}
    G_t = (V,E_t)
    \label{eq:graph}
\end{equation}

where:

\begin{itemize}
    \item $V$ is the set of locations or intersections.
    \item $E_t$ is the set of roads available at time $t$.
    \item $G_t$ represents the current road network.
\end{itemize}

The subscript $t$ indicates that the graph can change during the
simulation.

If a road is closed, it is removed from the current graph:

\begin{equation}
    e \notin E_t
    \qquad
    \text{if road } e \text{ is closed at time } t.
    \label{eq:closedroad}
\end{equation}

This prevents routing algorithms from selecting a closed road.





\subsection{Dynamic Road-Cost Model}

The travel cost of a road depends on its normal travel time,
congestion, weather, and hazards.

The proposed cost function is:

\begin{equation}
    C_e(t)
    =
    T_e
    \left(1+\alpha\rho_e(t)\right)
    W_e(t)
    +
    \beta H_e(t)
    \label{eq:roadcost}
\end{equation}

where:

\begin{itemize}
    \item $C_e(t)$ is the current travel cost of road $e$.
    \item $T_e$ is the normal travel time of road $e$.
    \item $\rho_e(t)$ is the congestion level.
    \item $\alpha$ controls the effect of congestion.
    \item $W_e(t)$ is the weather multiplier.
    \item $H_e(t)$ represents the hazard level.
    \item $\beta$ controls the hazard penalty.
\end{itemize}

A higher congestion or hazard level therefore increases the cost of
using the affected road.





\subsection{Request Arrival Model}

Requests can be generated using a Poisson process. The probability
of receiving $k$ requests during a time interval $t$ is:

\begin{equation}
    P(N(t)=k)
    =
    \frac{(\lambda t)^k e^{-\lambda t}}{k!}
    \label{eq:poisson}
\end{equation}

where:

\begin{itemize}
    \item $N(t)$ is the number of requests during the interval.
    \item $k$ is the number of requests.
    \item $\lambda$ is the average request arrival rate.
    \item $t$ is the duration of the interval.
\end{itemize}

This allows the simulation to reproduce both ordinary demand and
bursty demand.





\subsection{A* Pathfinding}

For every candidate node $n$, A* calculates:

\begin{equation}
    f(n)=g(n)+h(n)
    \label{eq:astar}
\end{equation}

where:

\begin{itemize}
    \item $g(n)$ is the actual cost from the starting node to $n$.
    \item $h(n)$ is the estimated cost from $n$ to the destination.
    \item $f(n)$ is the estimated total cost through $n$.
\end{itemize}

The node with the lowest relevant evaluation value is considered
first during the search.

The edge costs used by A* come from the dynamic road-cost model in
Equation~\eqref{eq:roadcost}.





\subsection{Estimated Time of Arrival}

For vehicle $i$ and request $j$, let $P_{ij}$ denote the selected
path.

The estimated travel time is:

\begin{equation}
    ETA_{ij}
    =
    \sum_{e\in P_{ij}} C_e(t)
    \label{eq:eta}
\end{equation}

where:

\begin{itemize}
    \item $ETA_{ij}$ is the estimated travel time from vehicle $i$
    to request $j$.
    \item $P_{ij}$ is the selected path.
    \item $C_e(t)$ is the current cost of road $e$.
\end{itemize}

The ETA values are used to construct the assignment matrix.





\subsection{Vehicle-Request Cost Matrix}

Suppose there are $n$ available vehicles and $m$ pending requests.
The resulting matrix is:

\begin{equation}
\mathbf{C}
=
\begin{bmatrix}
C_{11} & C_{12} & \cdots & C_{1m} \\
C_{21} & C_{22} & \cdots & C_{2m} \\
\vdots & \vdots & \ddots & \vdots \\
C_{n1} & C_{n2} & \cdots & C_{nm}
\end{bmatrix}
\label{eq:costmatrix}
\end{equation}

Here, $C_{ij}$ represents the cost of assigning vehicle $i$ to
request $j$.

For a basic model:

\begin{equation}
    C_{ij}=ETA_{ij}
    \label{eq:basiccost}
\end{equation}

The cost can later include additional penalties such as deadline
violations.





\subsection{Hungarian Assignment Model}

Let $x_{ij}$ be a binary decision variable:

\begin{equation}
x_{ij}
=
\begin{cases}
1, & \text{if vehicle } i \text{ is assigned to request } j,\\
0, & \text{otherwise.}
\end{cases}
\label{eq:assignment}
\end{equation}

The objective is to minimize total assignment cost:

\begin{equation}
    \min
    \sum_{i=1}^{n}
    \sum_{j=1}^{m}
    C_{ij}x_{ij}
    \label{eq:hungarianobjective}
\end{equation}

Each vehicle can be assigned to at most one simultaneous request:

\begin{equation}
    \sum_{j=1}^{m}x_{ij}
    \leq 1,
    \qquad
    \forall i
    \label{eq:vehicleconstraint}
\end{equation}

Each request can be assigned to at most one vehicle:

\begin{equation}
    \sum_{i=1}^{n}x_{ij}
    \leq 1,
    \qquad
    \forall j
    \label{eq:requestconstraint}
\end{equation}

The Hungarian algorithm is used to solve this assignment problem
efficiently.





\subsection{Deadline-Aware Assignment}

Let $D_j$ be the required deadline of request $j$. If the current
simulation time is $t$, the expected arrival time is:

\begin{equation}
    A_{ij}
    =
    t+ETA_{ij}
    \label{eq:arrival}
\end{equation}

The amount of lateness is:

\begin{equation}
    L_{ij}
    =
    \max
    \left(
        0,
        A_{ij}-D_j
    \right)
    \label{eq:lateness}
\end{equation}

If the vehicle arrives before the deadline, the lateness is zero.

The assignment cost can then be modified to:

\begin{equation}
    C_{ij}
    =
    ETA_{ij}
    +
    \mu L_{ij}
    \label{eq:deadlinecost}
\end{equation}

where $\mu$ controls the importance of meeting the deadline.

This allows the assignment stage to consider both travel time and
SLA compliance.





\subsection{Vehicle Routing with Time Windows}

For logistics missions involving multiple destinations, the project
uses the Vehicle Routing Problem with Time Windows (VRPTW).

Let $x_{ij}$ indicate whether a route travels directly from location
$i$ to location $j$.

The basic objective is:

\begin{equation}
    \min
    \sum_{i\in V}
    \sum_{j\in V}
    d_{ij}x_{ij}
    \label{eq:vrptw}
\end{equation}

where $d_{ij}$ is the travel cost between locations $i$ and $j$.

Each location may also have an acceptable service window:

\begin{equation}
    a_i
    \leq
    T_i
    \leq
    b_i
    \label{eq:timewindow}
\end{equation}

where:

\begin{itemize}
    \item $a_i$ is the earliest acceptable arrival time.
    \item $b_i$ is the latest acceptable arrival time.
    \item $T_i$ is the actual arrival time.
\end{itemize}





\subsection{Genetic Algorithm Fitness}

A candidate logistics route $R$ is evaluated using a fitness cost:

\begin{equation}
    F(R)
    =
    D(R)
    +
    \gamma P(R)
    \label{eq:fitness}
\end{equation}

where:

\begin{itemize}
    \item $D(R)$ is the travel cost of route $R$.
    \item $P(R)$ is the penalty caused by constraint violations.
    \item $\gamma$ controls the importance of those violations.
\end{itemize}

The Genetic Algorithm attempts to minimize $F(R)$ through selection,
crossover, and mutation.





\subsection{Reinforcement-Learning Model}

The idle-vehicle repositioning problem is represented as a Markov
Decision Process:

\begin{equation}
    \mathcal{M}
    =
    (\mathcal{S},\mathcal{A},P,R,\gamma)
    \label{eq:mdp}
\end{equation}

where:

\begin{itemize}
    \item $\mathcal{S}$ is the set of possible system states.
    \item $\mathcal{A}$ is the set of possible repositioning actions.
    \item $P$ represents state-transition probabilities.
    \item $R$ represents the reward function.
    \item $\gamma$ is the discount factor.
\end{itemize}

The state can include information such as vehicle locations,
vehicle availability, demand distribution, and road conditions.





\subsection{Q-Learning Update}

The Q-value is updated using:

\begin{equation}
\begin{aligned}
Q(s,a)
\leftarrow
Q(s,a)
+
\eta
\Big[
r
+
\gamma
\max_{a'}
Q(s',a')
-
Q(s,a)
\Big]
\end{aligned}
\label{eq:qlearning}
\end{equation}

where:

\begin{itemize}
    \item $s$ is the current state.
    \item $a$ is the selected action.
    \item $r$ is the reward received.
    \item $s'$ is the next state.
    \item $\eta$ is the learning rate.
    \item $\gamma$ is the discount factor.
    \item $a'$ represents a possible action in the next state.
\end{itemize}

The policy is trained to learn which repositioning actions are useful
for future demand.





\subsection{Repositioning Reward}

The repositioning policy must balance faster future response with
the additional cost of moving idle vehicles.

A simplified reward function is:

\begin{equation}
    r
    =
    -w_TT
    -
    w_EE
    +
    w_SS
    \label{eq:reward}
\end{equation}

where:

\begin{itemize}
    \item $T$ represents response time.
    \item $E$ represents additional energy consumption.
    \item $S$ represents service performance.
    \item $w_T$ is the response-time weight.
    \item $w_E$ is the energy-consumption weight.
    \item $w_S$ is the service-performance weight.
\end{itemize}

The reward encourages the system to reduce response time and improve
service while avoiding unnecessary repositioning.





\subsection{Deep Q-Network Extension}

If the state space becomes too large for a Q-table, a Deep Q-Network
can be used to approximate the Q-value:

\begin{equation}
    Q(s,a;\theta)
    \label{eq:dqn}
\end{equation}

where $\theta$ represents the trainable parameters of the neural
network.

DQN is considered an extension of the initial Q-Learning
implementation rather than a requirement for the basic simulation.





\subsection{Vehicle Battery Model}

For electric vehicles, the battery state is normalized between zero
and one:

\begin{equation}
    0
    \leq
    B_i(t)
    \leq
    1
    \label{eq:battery}
\end{equation}

After movement, the battery state can be updated as:

\begin{equation}
    B_i(t+\Delta t)
    =
    B_i(t)
    -
    E_i(\Delta t)
    \label{eq:batteryupdate}
\end{equation}

where:

\begin{itemize}
    \item $B_i(t)$ is the battery level of vehicle $i$.
    \item $E_i(\Delta t)$ is the energy consumed during the interval.
    \item $\Delta t$ is the simulation time interval.
\end{itemize}

A vehicle with insufficient battery should not be selected for a
mission that it cannot complete safely.





\section{System Architecture}

AetherGrid follows a layered architecture so that the simulation
engine, algorithms, and user interface remain separated.

\subsection{Simulation Core}

The simulation core contains:

\begin{itemize}
    \item the road or grid network,
    \item vehicle agents,
    \item request agents,
    \item environmental conditions,
    \item vehicle battery state,
    \item the seeded request generator, and
    \item the discrete-event or tick-based simulation loop.
\end{itemize}

The core must be able to run without the web interface. This provides
a headless mode for testing, benchmarking, and continuous integration.

\subsection{Algorithm Layer}

The algorithm layer contains independent strategy modules:

\begin{itemize}
    \item Dynamic A*,
    \item Dijkstra,
    \item HPA*,
    \item Hungarian assignment,
    \item Genetic Algorithm based VRPTW,
    \item Q-Learning, and
    \item optional DQN.
\end{itemize}

Each strategy should have a clear interface so that different
strategies can be compared without changing the simulation engine.

\subsection{Web and Visualization Layer}

The web layer provides Mission Control functionality.

The dashboard will display:

\begin{itemize}
    \item the current road network,
    \item vehicle locations,
    \item active requests,
    \item assigned routes,
    \item road conditions,
    \item vehicle status,
    \item battery information,
    \item response-time statistics,
    \item SLA compliance, and
    \item simulation events.
\end{itemize}

The proposed interface uses WebSocket-based state updates. The
backend can publish state changes at approximately $10$ Hz while the
browser renders the visual scene at up to $60$ frames per second.

\subsection{Scenario Laboratory}

The dashboard will also contain a scenario laboratory through which
the user can inject events such as:

\begin{itemize}
    \item road closures,
    \item congestion,
    \item hazards,
    \item weather changes,
    \item demand bursts, and
    \item new emergency requests.
\end{itemize}

This makes it possible to demonstrate how the dispatch strategies
react to controlled disruptions.





\section{Proposed Technology Stack}

The following technology stack is proposed for implementation.

\begin{table}[ht]
\centering
\renewcommand{\arraystretch}{1.35}
\begin{tabularx}{\textwidth}{>{\bfseries}p{4.0cm}X}
\toprule
Layer & Proposed Technology / Approach \\
\midrule
Programming Language & Python 3.11 \\
Numerical Computation & NumPy and standard Python libraries \\
Graph Algorithms & Custom implementations of A*, Dijkstra, and HPA* \\
Assignment & Hungarian algorithm implementation \\
Route Optimization & Genetic Algorithm for VRPTW \\
Reinforcement Learning & Q-Learning with optional DQN extension \\
Simulation & Seeded discrete-event / tick-based simulation engine \\
Backend & Python web backend with WebSocket communication \\
Frontend & HTML, CSS, JavaScript, and Canvas-based visualization \\
Testing & Unit tests and property-based tests \\
Version Control & Git and GitHub \\
Continuous Integration & Automated seed-stable regression and invariant tests \\
\bottomrule
\end{tabularx}
\caption{Proposed technology stack for AetherGrid AI.}
\label{tab:techstack}
\end{table}

The project intentionally uses established algorithms and standard
software components. The primary engineering challenge is their
correct integration into one reproducible system.





\section{Implementation Methodology}

The project will be implemented incrementally so that each stage
produces a working and testable result.

\subsection{Phase 1: Simulation Foundation}

The first phase will implement:

\begin{itemize}
    \item the road graph,
    \item vehicle and request models,
    \item environmental conditions,
    \item the battery state model,
    \item the seeded request generator, and
    \item the simulation loop.
\end{itemize}

\textbf{Exit criterion:} a seed-stable headless simulation run.

\subsection{Phase 2: Routing and Assignment}

The second phase will implement:

\begin{itemize}
    \item Dynamic A*,
    \item Dijkstra,
    \item HPA*,
    \item ETA matrix generation,
    \item Hungarian assignment,
    \item deadline-aware costs, and
    \item baseline greedy dispatch.
\end{itemize}

\textbf{Exit criterion:} routing and assignment tests pass against
known small examples and baseline strategies.

\subsection{Phase 3: Logistics and Reinforcement Learning}

The third phase will implement:

\begin{itemize}
    \item VRPTW representation,
    \item Genetic Algorithm route optimization,
    \item Q-Learning for idle-vehicle repositioning, and
    \item optional DQN extension.
\end{itemize}

\textbf{Exit criterion:} each strategy can be run independently and
compared under identical seeded scenarios.

\subsection{Phase 4: Mission Control Dashboard}

The fourth phase will implement:

\begin{itemize}
    \item live map visualization,
    \item vehicle telemetry,
    \item request monitoring,
    \item route visualization,
    \item scenario injection,
    \item strategy switching, and
    \item performance charts.
\end{itemize}

\textbf{Exit criterion:} a complete scenario can be run from the
dashboard while the dispatch strategy is changed during operation.

\subsection{Phase 5: Testing and Evaluation}

The final phase will focus on:

\begin{itemize}
    \item performance benchmarking,
    \item reproducibility,
    \item stress testing,
    \item disruption testing,
    \item baseline comparison, and
    \item final documentation.
\end{itemize}

\textbf{Exit criterion:} all predefined evaluation experiments have
been executed and documented.





\section{Testing and Validation}

AetherGrid will use both functional tests and performance tests.

\subsection{Functional Invariants}

Important system properties include:

\begin{itemize}
    \item No route may traverse a closed road.
    \item A vehicle cannot be assigned to two simultaneous requests.
    \item A request cannot be assigned to two vehicles.
    \item Battery levels remain within their valid range.
    \item Completed vehicles return to the available pool.
    \item The same random seed produces the same simulation trace.
\end{itemize}

\subsection{Assignment Validation}

For small problem sizes, the Hungarian solution can be compared
against brute-force enumeration. For larger problem sizes, it can
be compared against a greedy nearest-vehicle baseline.

The expected property is that the optimized assignment should not
have a higher total assignment cost than the corresponding greedy
assignment under the same cost matrix.

\subsection{Routing Validation}

A* and Dijkstra will be tested on small graphs for which the optimal
path cost can be calculated independently.

The tests will also verify that introducing a road closure changes
the route when the original route is no longer valid.

\subsection{Reinforcement-Learning Validation}

The RL component will be evaluated by comparing:

\begin{enumerate}
    \item no repositioning,
    \item random repositioning,
    \item heuristic repositioning, and
    \item learned repositioning.
\end{enumerate}

The comparison will use identical demand scenarios and random seeds.





\section{Evaluation Metrics}

The following metrics will be used to compare the proposed
strategies.

\subsection{Average Response Time}

The average response time is:

\begin{equation}
    \overline{T}
    =
    \frac{1}{N}
    \sum_{i=1}^{N}T_i
    \label{eq:avgresponse}
\end{equation}

where $N$ is the number of completed requests and $T_i$ is the
response time of request $i$.

A lower value indicates faster average response.

\subsection{SLA Compliance Rate}

The percentage of requests completed within their deadlines is:

\begin{equation}
    SLA_{\mathrm{rate}}
    =
    \frac{N_{\mathrm{on\text{-}time}}}
    {N_{\mathrm{total}}}
    \times 100
    \label{eq:sla}
\end{equation}

where:

\begin{itemize}
    \item $N_{\mathrm{on\text{-}time}}$ is the number of requests
    completed within their deadlines.
    \item $N_{\mathrm{total}}$ is the total number of requests.
\end{itemize}

A higher value represents better deadline compliance.

\subsection{Average Travel Distance}

The average distance travelled per mission is:

\begin{equation}
    \overline{D}
    =
    \frac{1}{N}
    \sum_{i=1}^{N}D_i
    \label{eq:avgdistance}
\end{equation}

where $D_i$ is the distance travelled during mission $i$.

\subsection{Average Energy Consumption}

The average energy consumption can be measured as:

\begin{equation}
    \overline{E}
    =
    \frac{1}{N}
    \sum_{i=1}^{N}E_i
    \label{eq:avgen}
\end{equation}

where $E_i$ is the energy consumed by mission $i$.

\subsection{Algorithm Runtime}

The runtime of a decision-making algorithm is:

\begin{equation}
    T_{\mathrm{runtime}}
    =
    T_{\mathrm{finish}}
    -
    T_{\mathrm{start}}
    \label{eq:runtime}
\end{equation}

The measured runtime will be compared with the five-second batch
interval defined in Equation~\eqref{eq:batch}.





\section{Overall Optimization Objective}

AetherGrid attempts to balance response time, service quality,
energy consumption, and travel distance.

A simplified multi-objective cost is:

\begin{equation}
    J
    =
    w_1\overline{T}
    +
    w_2(1-SLA_{\mathrm{rate}})
    +
    w_3\overline{E}
    +
    w_4\overline{D}
    \label{eq:overall}
\end{equation}

where:

\begin{itemize}
    \item $\overline{T}$ is the average response time.
    \item $SLA_{\mathrm{rate}}$ is the SLA compliance rate.
    \item $\overline{E}$ is the average energy consumption.
    \item $\overline{D}$ is the average travel distance.
    \item $w_1,w_2,w_3,w_4$ are weights representing the importance
    of the corresponding objectives.
\end{itemize}

The system seeks to minimize $J$ while satisfying the operational
constraints of the simulation.

The exact weights will be selected during experimentation rather than
being presented as predetermined results.





\section{Experimental Design}

The evaluation will use controlled and reproducible scenarios.

\subsection{Baseline Strategies}

The intelligent strategies will be compared with simpler methods:

\begin{itemize}
    \item \textbf{Nearest Vehicle:} assigns the closest available
    vehicle to each request.

    \item \textbf{First-Come-First-Served:} processes requests in
    arrival order.

    \item \textbf{Greedy Assignment:} repeatedly selects the
    currently cheapest vehicle-request pair.

    \item \textbf{Hungarian Assignment:} performs batch-level global
    assignment.

    \item \textbf{Hungarian + RL:} combines optimized dispatch with
    learned idle-vehicle repositioning.
\end{itemize}

\subsection{Scenario Categories}

The following scenarios will be tested:

\begin{enumerate}
    \item Normal traffic and normal request volume.
    \item High traffic.
    \item Bursty emergency demand.
    \item Multiple simultaneous requests.
    \item Road closure during an active simulation.
    \item Hazard introduced during an active route.
    \item Uneven spatial demand.
    \item High vehicle utilization.
    \item Mixed emergency and logistics requests.
\end{enumerate}

Each strategy will be evaluated on identical seeded scenarios to
ensure that the comparison is fair.





\section{Evaluation Hypotheses}

The project will use measurable hypotheses rather than claiming
performance improvements before the experiments are completed.

\subsection{Hypothesis 1: Batch Matching}

Under bursty demand, the hypothesis is that global batch assignment
will perform better than a purely greedy nearest-vehicle strategy.

The comparison will use identical request streams and vehicle states.

The hypothesis will be considered unsupported if the measured
difference is negligible across the selected scenarios.

\subsection{Hypothesis 2: Predictive Repositioning}

The second hypothesis is that repositioning idle vehicles toward
anticipated demand can reduce future response time compared with
dispatch without repositioning.

The additional energy and movement cost will also be measured. An
improvement in response time alone will not be considered sufficient
if the additional operational cost becomes excessive.

\subsection{Hypothesis 3: Disruption Handling}

The third hypothesis is that the dynamic cost field will allow the
system to respond to hazards and road closures without violating
routing constraints.

The test will verify that:

\begin{itemize}
    \item affected routes are recalculated,
    \item closed roads are not traversed,
    \item vehicles continue operating when alternative routes exist,
    \item no invalid vehicle assignment is produced.
\end{itemize}





\section{Reproducibility and Continuous Integration}

Reproducibility is an important design requirement.

Every simulation run will use an explicit random seed. When the same
configuration and seed are used, the generated request sequence and
simulation behavior should be reproducible.

Automated tests will be executed whenever changes are made to the
project.

Important automated checks include:

\begin{itemize}
    \item no closed-edge traversal,
    \item valid battery bounds,
    \item valid vehicle-request assignments,
    \item deterministic seeded runs,
    \item assignment-cost comparison against baseline methods,
    \item route validity, and
    \item regression testing of core simulation behavior.
\end{itemize}

This approach allows the simulation engine to operate in a headless
mode for automated testing while the same engine can drive the live
dashboard for demonstrations.





\section{Scalability}

The architecture is designed so that the simulation engine is
separated from the visualization layer.

The assignment problem can be solved using the Hungarian algorithm
with cubic time complexity:

\begin{equation}
    O(n^3)
    \label{eq:hungariancomplexity}
\end{equation}

for an $n\times n$ assignment matrix.

This makes the assignment stage practical for moderate fleet sizes
within the fixed dispatch interval. For larger graphs, HPA* can be
used to reduce the routing search space.

The separation between the simulation engine and dashboard also
ensures that rendering workload does not directly determine the
correctness of the simulation.





\section{Project Scope}

\subsection{Included in Scope}

The project will include:

\begin{itemize}
    \item dynamic road-network simulation,
    \item heterogeneous vehicle simulation,
    \item emergency and logistics requests,
    \item congestion and environmental effects,
    \item road closures and hazards,
    \item Dynamic A*,
    \item Dijkstra,
    \item HPA*,
    \item Hungarian assignment,
    \item Genetic Algorithm based VRPTW,
    \item Q-Learning,
    \item optional DQN,
    \item vehicle battery modeling,
    \item real-time dashboard,
    \item scenario laboratory,
    \item reproducible experiments, and
    \item baseline performance comparison.
\end{itemize}

\subsection{Outside the Initial Scope}

The following are not required for the first complete version:

\begin{itemize}
    \item deployment on a real city-wide transportation network,
    \item direct control of real vehicles,
    \item dependence on proprietary real-time traffic data,
    \item production-grade autonomous driving,
    \item guaranteed real-world emergency response performance.
\end{itemize}

The system is a simulation and engineering demonstration. Its
experimental results should therefore be interpreted within the
conditions of the implemented simulation.





\section{Risks and Mitigation}

\subsection{Scope Creep}

The project combines several algorithm families. To control the
scope, every algorithm will be implemented behind a fixed interface
and developed in phases.

\subsection{Real-Time Performance}

If routing becomes expensive on larger graphs, HPA* can be introduced
for long-haul searches. Assignment runtime will also be measured
against the five-second batch interval from the beginning of
development.

\subsection{Reinforcement-Learning Instability}

The RL component will remain advisory and will only control idle
vehicle repositioning. It will never override an active dispatch.
Therefore, an incompletely trained policy cannot invalidate an
emergency assignment.

\subsection{Evaluation Ambiguity}

The experiments will be defined before final benchmarking. All
strategies will be evaluated using identical seeds, scenarios, and
metrics.





\section{Expected Outcomes}

At the completion of the project, the following outcomes are
expected:

\begin{enumerate}
    \item A working seeded urban dispatch simulation engine.

    \item A dynamic road graph capable of representing congestion,
    weather, hazards, and road closures.

    \item A routing system supporting Dynamic A*, Dijkstra, and HPA*.

    \item A batch assignment system using the Hungarian algorithm.

    \item A logistics routing module based on Genetic Algorithm and
    VRPTW.

    \item A reinforcement-learning module for idle-vehicle
    repositioning.

    \item A real-time Mission Control dashboard.

    \item A scenario laboratory for controlled disruptions.

    \item Automated tests for important system invariants.

    \item Reproducible benchmark results comparing multiple dispatch
    strategies.
\end{enumerate}

The final result will therefore be a complete software system rather
than an isolated implementation of a single algorithm.





\section{Project Deliverables}

The planned deliverables are:

\begin{table}[ht]
\centering
\renewcommand{\arraystretch}{1.35}
\begin{tabularx}{\textwidth}{>{\bfseries}p{3.5cm}X}
\toprule
Deliverable & Description \\
\midrule
Simulation Engine & Seeded urban road, vehicle, request, and environment simulator. \\
Routing Module & Dynamic A*, Dijkstra, and HPA* implementations. \\
Assignment Module & Hungarian vehicle-request matching with deadline-aware costs. \\
Logistics Module & Genetic Algorithm based VRPTW route optimization. \\
RL Module & Q-Learning based idle-vehicle repositioning with optional DQN extension. \\
Mission Control & Real-time browser dashboard for visualization and monitoring. \\
Scenario Lab & Controlled injection of traffic, hazards, closures, and demand bursts. \\
Testing Suite & Unit, property, invariant, and reproducibility tests. \\
Benchmark Suite & Seeded comparison against greedy and other baseline strategies. \\
Documentation & Technical report, architecture documentation, and usage instructions. \\
\bottomrule
\end{tabularx}
\caption{Planned AetherGrid deliverables.}
\label{tab:deliverables}
\end{table}





\section{Conclusion}

AetherGrid AI proposes a practical simulation-based approach to
dynamic urban dispatch. Instead of treating routing, assignment, and
future vehicle positioning as one problem, the system separates them
into cooperating layers.

The dynamic road model represents changing travel conditions. The
pathfinding layer calculates routes using current road costs. The
Hungarian algorithm provides global batch-level assignment for
one-to-one dispatch. The Genetic Algorithm handles multi-stop
logistics routing with time windows. Finally, the reinforcement-
learning layer provides limited predictive behavior by repositioning
idle vehicles between dispatch cycles.

The project is intentionally designed around established algorithms
rather than claiming algorithmic novelty. Its contribution is the
engineering integration of these methods into a reproducible,
testable, real-time simulation environment.

The final system will allow different strategies to be tested under
the same controlled scenarios. This makes it possible to determine,
using measurable evidence, whether global assignment and predictive
idle-vehicle repositioning actually improve response performance,
service compliance, and operational efficiency.

\end{document}
